\scanpage{88}
\mspage{5}
% archive sheet number: 87
morphismes de transition, \inserted{entre ces faisceaux, \unsure{identifiés comme} le faisceau $\mu_{\ell^\nu}$ sur $X_0$} On trouve que si $S_{\deleted{\unsure{\beta}}\alpha}$ est de degré $n_\alpha$ sur $S$, donc $S_\beta$ de degré $\inserted{n_{\alpha\beta}=}n_\beta/n_\alpha$ sur $S_\alpha$ (si $\beta\geq\alpha$), alors le morphisme de transition est la multiplication par $n_{\alpha\beta}$. \inserted{(On est ramené à faire \unsure{cette} vérification pour $X=S$\dots)} On en conclut aussitôt que $\underline{\unsure{\Phi|\varphi}}^2_\nu=\varinjlim(\mu_{\ell^\nu})$ \deleted{\unsure{est isomorphe}}\inserted{est également nul}. \deleted{\unsure{plus simplement}} \uline{Donc} $\underline\Phi^i_\nu=0$ \uline{pour \unsure{tt|$\forall$} $i,\nu$}, \deleted{d'où} \deleted{grâce au fait} On en déduit grâce à (1.14) que $\Phi^i_\nu=0$ pour tt $i,\nu$, puis, grâce à (1.9) que
\[
H^i(X_s,\Lambda_\nu)\xrightarrow{\sim}H^i(X_{\bar{\unsure{\eta|y}}},\Lambda_\nu)
\tag*{(1.16.1)}
\]
($X$ lisse sur $S$, \deleted{\unsure{supposant}} $\ell\ne p$) : \deleted{c'est}\inserted{on \deleted{\illegible{1}} \unsure{retrouve}} le \uline{th. de spécialisation} de la cohomologie pour un morphisme propre et lisse. (SGA 4 $\overline{\underline{\mathrm{XVI}}}$ \dots). [Bien entendu, la \unsure{même} démonstration marcherait avec des coefficients constants tordus).

Compte tenu de la nature locale des $\underline\Phi^i_\nu$ en général, on conclut de ce qui précède la

\uline{Proposition} 1.17. Soit $T$ l'ensemble des $x\in X_0$ en lesquels $f\inserted{:X\longrightarrow S}$ n'est pas lisse, \inserted{et supposons $\ell\ne p$.} Alors les faisceaux $\underline\Phi^i_\nu$ des cycles évanescents sont à support $\subset T$.

\deleted{les points de $T$ étant \unsure{fermés},}\deleted{\unsure{sur}}

1.18. Supposons p.ex. que $T$ soit \uline{discret}. Alors dans \deleted{$\unsure{E_2^{pq}}$} (1.14) on a $E_2^{pq}=0$ pour $p\ne0$, donc on trouve des isomorphismes
